% =====================================================================
% Clase -- Geometría 7° -- Trimestre I -- Semana 03
% Sesión 003 (mar 15 sep., 6:55, 55 min)
% Tema: Ángulos y giros
% Hilo conductor del trimestre I: «Escher: el arte de mover figuras»
% Tema 2 de la guía: Ángulos y giros (tema:giros)
% Material del docente (no se entrega a los estudiantes).
% ---------------------------------------------------------------------
% BORRADOR: la fila 003 de curriculo/programacion.csv está VACÍA. El tema, el subtema y
% el DBA de abajo se toman de la propuesta del trimestre I de curriculo/plan-anual.md,
% que la docente aún no ha aprobado, y NO se han escrito en programacion.csv.
% La propuesta marca tarea para esta sesión: no se escribió porque el banco no tiene
% ejercicios verificados de este tema para 7° (ver PENDIENTES.md).
% =====================================================================
\documentclass[11pt]{article}
\makeatletter\def\input@path{{../../../../../plantillas/estilo/}}\makeatother
\usepackage{mmcantillo}
\guiadelestudiante{../../guia-didactica/guia-periodo-I-transformaciones}

\tipodocumento{Clase}
\asignatura{Geometría}
\titulo{Semana 3: el ángulo como giro}
\grado{7°}
\periodo{I}

% DBA: PROPUESTO, sin aprobar — matematicas grado 7 · DBA 5 — «Observa objetos
% tridimensionales desde diferentes puntos de vista, los representa según su ubicación y
% los reconoce cuando se transforman mediante rotaciones, traslaciones y reflexiones.»
% Esta sesión prepara la rotación del DBA 5: el ángulo de giro y su sentido.

\begin{document}

\encabezado*

\begin{center}
{\renewcommand{\arraystretch}{1.3}
\begin{tabularx}{\linewidth}{@{}l l l X l@{}}
  \toprule
  \textbf{Sesión} & \textbf{Fecha} & \textbf{Min} & \textbf{Subtema} & \textbf{Entregables}\\
  \midrule
  003 & mar 15 sep., 6:55 & 55 & El ángulo como giro: sentido horario y antihorario,
      medición con transportador, giros de $90^\circ$, $180^\circ$, $270^\circ$ y
      $360^\circ$ & --- (ver nota)\\
  \bottomrule
\end{tabularx}}
\end{center}

La semana 02 clasificó los ángulos por su medida: agudo, recto, obtuso, llano y completo.
Esta sesión les da un segundo significado, el que hace falta para todo el trimestre: el
ángulo deja de ser una abertura quieta y pasa a medir \emph{cuánto giró} una figura. Con
eso queda abierta la puerta a las rotaciones, las traslaciones y las reflexiones de
Escher. Referencia: \refguia{tema:giros}.

\begin{nota}
\textbf{Pendientes de esta sesión.} (1) La fila 003 de \texttt{programacion.csv} está
vacía: el tema y el DBA de arriba salen de la propuesta del plan anual, que falta
aprobar. (2) La propuesta marca \emph{tarea} para hoy; no se escribió porque el banco no
tiene ejercicios verificados de giros para 7° (solo \texttt{angulos-7-003}, que ya usa la
guía). La práctica de la sesión va en el cuaderno y se revisa en clase.
\end{nota}

% ---------------------------------------------------------------------
\seccion{Sesión 003 — martes 15 de septiembre · 6:55 · 55 min}

\textbf{Objetivo.} Describir un giro con su ángulo y su sentido, medir ángulos de giro con
el transportador y reconocer que girar $a$ grados en un sentido equivale a girar
$360^\circ - a$ en el contrario.

\textbf{Materiales.} Tablero; transportador y regla (cada estudiante); una escuadra o un
recorte de cartulina en forma de «L» por pareja; un reloj de manecillas dibujado en el
tablero; la guía del estudiante.

\begin{center}
{\renewcommand{\arraystretch}{1.3}
\begin{tabularx}{\linewidth}{@{}l l X@{}}
  \toprule
  \textbf{Momento} & \textbf{Min} & \textbf{Actividad}\\
  \midrule
  Enganche & 0--8 & Leer el recuadro del hilo del Tema 2: la estrella de la Alhambra que
    al girarla vuelve a verse igual. Preguntar: «¿cuánto hay que girar una estrella de
    ocho puntas para que se vea igual?» Dejar que tanteen; la respuesta ($45^\circ$)
    vuelve al final. Contar en una frase quién fue M. C. Escher: un artista holandés que
    llenó sus grabados de figuras que se repiten girando y deslizándose, y que sacó la
    idea de los mosaicos de la Alhambra, en Granada.\\
  Del ángulo quieto al giro & 8--18 & Recordar la clasificación de la semana 02 con el
    transportador. Luego dibujar dos rayos con el mismo vértice y preguntar: «¿este
    ángulo es una esquina o un movimiento?» Definir el giro: \textbf{centro} (el punto
    fijo), \textbf{ángulo} y \textbf{sentido}. Antihorario $=$ positivo, horario $=$
    negativo. Mostrarlo girando la escuadra sobre el puesto.\\
  Con el reloj & 18--28 & El minutero da una vuelta en 60 min, así que gira
    $360^\circ \div 60 = 6^\circ$ por minuto. Preguntas al vuelo: ¿cuánto gira en 5
    minutos ($30^\circ$)?, ¿en 15 ($90^\circ$)?, ¿en media hora ($180^\circ$)? ¿En qué
    sentido? (Horario: por eso se llama así.) Cerrar con la rueda de la fortuna de la
    guía: 12 cabinas, $30^\circ$ de cabina a cabina.\\
  Giros equivalentes & 28--40 & Cada pareja marca una esquina de su escuadra con un punto
    de color, la apoya sobre el cuaderno y la gira: $90^\circ$, $180^\circ$, $270^\circ$ y
    $360^\circ$ antihorario, dibujando la posición cada vez. Después, $90^\circ$ en
    sentido horario. Comparar con el giro de $270^\circ$ antihorario: caen en el mismo
    lugar. Escribir la regla en el tablero: girar $a$ en un sentido $=$ girar
    $360^\circ - a$ en el contrario, porque juntos dan una vuelta completa.\\
  Práctica & 40--52 & En el cuaderno, individual (ver las respuestas abajo).\\
  Cierre & 52--55 & Volver a la estrella de ocho puntas: $360^\circ \div 8 = 45^\circ$.
    Idea clave: \emph{un giro se dice con tres datos — centro, ángulo y sentido}. Anunciar
    que la próxima semana la figura no gira: se desliza.\\
  \bottomrule
\end{tabularx}}
\end{center}

\textbf{Práctica de la sesión (en el cuaderno) y sus respuestas.}
\begin{enumerate}
  \item ¿A qué giro en sentido horario equivale cada uno de estos giros antihorarios:
    $90^\circ$, $120^\circ$, $210^\circ$, $300^\circ$?
    \\\emph{Respuesta:} $270^\circ$, $240^\circ$, $150^\circ$ y $60^\circ$
    (en cada caso, $360^\circ$ menos el ángulo dado).
  \item El minutero de un reloj marca las 12:00. ¿Qué hora marca después de un giro de
    $180^\circ$? ¿Y después de dos giros de $270^\circ$?
    \\\emph{Respuesta:} 12:30. Dos giros de $270^\circ$ son $540^\circ = 360^\circ +
    180^\circ$: una vuelta completa y media vuelta más, así que otra vez las 12:30.
  \item Una hélice de $3$ aspas iguales se ve igual cada vez que gira cierto ángulo.
    ¿Cuál es el menor giro que la deja igual? ¿Y si tuviera $6$ aspas?
    \\\emph{Respuesta:} $360^\circ \div 3 = 120^\circ$; con $6$ aspas,
    $360^\circ \div 6 = 60^\circ$.
  % Ejercicio: angulos-7-003
  \item (Del Tema 2 de la guía.) Laura dice que girar una figura $270^\circ$ en sentido
    antihorario la deja igual que girarla $90^\circ$ en sentido horario. ¿Tiene razón?
    \\\emph{Respuesta:} Sí. Los dos giros suman $270^\circ + 90^\circ = 360^\circ$, una
    vuelta completa, así que llevan la figura al mismo lugar. Se comprueba girando la
    escuadra.
\end{enumerate}

\begin{nota}
\textbf{Error típico.} Confundir el ángulo de giro con la posición final: dos figuras
pueden acabar en el mismo sitio con giros distintos ($90^\circ$ horario y $270^\circ$
antihorario). Conviene exigir siempre las dos palabras juntas —el número y el sentido—
cuando un estudiante describa un giro.

\textbf{Sobre las imágenes de Escher.} Sus grabados tienen derechos de autor: se describen
en voz alta o se buscan en clase, pero no se reproducen en el material impreso.
\end{nota}

\end{document}
